\documentclass[11pt,a4paper]{article}

\usepackage[utf8]{inputenc}
\usepackage[T1]{fontenc}
\usepackage[spanish,es-tabla]{babel}
\usepackage{amsmath,amssymb,amsthm,mathtools}
\usepackage{booktabs}
\usepackage{enumitem}
\usepackage{geometry}
\usepackage{setspace}
\usepackage{hyperref}
\usepackage{microtype}

\geometry{margin=1in}
\onehalfspacing

\hypersetup{
  pdftitle={Nota formal: Sinergia irreducible, Res, excess3 y PID},
  pdfauthor={Johel Padilla-Villanueva},
  pdfkeywords={RECD, Res, excess3, PID, O-information, sinergia irreducible}
}

\theoremstyle{definition}
\newtheorem{definition}{Definici\'on}[section]
\newtheorem{remark}{Observaci\'on}[section]
\newtheorem{example}{Ejemplo}[section]
\theoremstyle{plain}
\newtheorem{proposition}{Proposici\'on}[section]
\newtheorem{lemma}{Lema}[section]
\newtheorem{theorem}{Teorema}[section]
\newtheorem{corollary}{Corolario}[section]
\newtheorem{conjecture}{Conjetura}[section]
\theoremstyle{remark}
\newtheorem{scholium}{Scholium}[section]

\newcommand{\Res}{\mathrm{Res}}
\newcommand{\TC}{\mathrm{TC}}
\newcommand{\Syn}{\mathrm{Syn}}
\newcommand{\Surp}{\mathrm{Surp}}
\newcommand{\excess}{\mathrm{excess3}}
\newcommand{\Rtwo}{\mathcal{R}_{2}}
\newcommand{\Rind}{\mathcal{R}_{\mathrm{ind}}}
\newcommand{\Rpair}{\mathcal{R}_{\mathrm{pair}}}
\newcommand{\Jwin}{\mathcal{J}}
\newcommand{\KL}{D_{\mathrm{KL}}}
\newcommand{\claim}{\mathsf{Claim}}
\newcommand{\RECD}{\mathrm{RECD}}

\title{\textbf{Nota formal}\\[0.4em]
\large Sinergia irreducible: residuo $\Res$, proxy $\excess$,\\
PID y condiciones de aparici\'on\\[0.6em]
\normalsize (Punto 2 del roadmap te\'orico; asume Nota formal 1a--1c)}

\author{
Johel Padilla-Villanueva\\
\textit{Department of Environmental Health, University of Puerto Rico}\\
\texttt{johelpadilla@gmail.com}
}

\date{2026-08-01 \quad$\cdot$\quad v0.1}

\begin{document}
\maketitle

\begin{abstract}
Asumiendo el canon de la Nota formal~1 (Nivel~3 $:=$ residuo
$\Res(\Jwin;\Rtwo)$), esta nota caracteriza matem\'aticamente ese residuo,
lo sit\'ua frente a multiinformaci\'on, O-information y PID, y
\emph{delimita} el estatus de $\excess=0.6\,\Syn+0.4\,\Surp$ como proxy.
Se demuestran identidades y desigualdades exactas (TC como KL a
independencia; monoton\'ia del residuo al enriquecer la clase de
reducci\'on; separaci\'on Syn/Surp), se exhiben contraejemplos que impiden
identificar $\excess$ con un \'atomo PID, y se formulan condiciones
necesarias/suficientes operativas para la aparici\'on de sinergia
irreducible, m\'as el comportamiento bajo transformaciones mon\'otonas y
proyecciones a pares. El resultado es una identidad te\'orica fuerte del
Nivel~3 \emph{sin} sobreclaim de PID completo.
\end{abstract}

\tableofcontents
\vspace{1em}

% ============================================================
\section{Dependencia de la Nota 1 y pureza}
\label{sec:dep}

\begin{scholium}[Canon heredado]
Valen sin reabrir: tres planos (claim / estimador / reloj);
A1--A3; Nivel~3 $:=$ $\Res(\Jwin;\Rtwo)>\theta_3$;
$\excess$ es proxy, no definici\'on; anidamiento Tipo~II;
M-I-1 (no se demuestra el acto de ser).
\end{scholium}

\begin{scholium}[Anti-overclaim]
Nada de lo siguiente afirma que $\excess$ sea un \'atomo axiom\'atico de
PID, ni causalidad, ni el reloj RECD completo (T-IV / M-IV-2).
\end{scholium}

Notaci\'on: $\Sigma=S_m$ (alfabeto Bandt--Pompe), $P\in\mathcal{P}(\Sigma^d)$
denota una ley joint (t\'ipicamente la emp\'irica $\Jwin_t^{(w)}$),
$P_i$ sus m\'argenes, $P_{ij}$ las bipolares.

% ============================================================
\section{Jerarqu\'ia de clases de reducci\'on}
\label{sec:classes}

\begin{definition}[Tres clases anidadas]
\label{def:classes}
\begin{align}
\Rind
&:=\Bigl\{Q\in\mathcal{P}(\Sigma^d):
Q(s)=\prod_{i=1}^d q_i(s_i)\Bigr\},\\
\Rpair
&:=\Bigl\{Q:
\log Q(s)
=\sum_i\theta_i(s_i)+\sum_{i<j}\theta_{ij}(s_i,s_j)
-\psi(\theta)\Bigr\}
\end{align}
(familia exponencial de interacciones de orden $\le 2$, con funci\'on
de partici\'on $\psi$), y $\Rtwo$ denota, en un protocolo dado, una de
$\Rind$, $\Rpair$, o una clase inter\-media declarada.
Siempre $\Rind\subset\Rpair\subset\mathcal{P}(\Sigma^d)$ (la independencia
es el caso $\theta_{ij}\equiv 0$).
\end{definition}

\begin{definition}[Residuos KL]
\label{def:ResKL}
Con divergencia de Kullback--Leibler $\KL(P\|Q)=\sum_s P(s)\log_2\frac{P(s)}{Q(s)}$
(en bits; $0\log 0:=0$),
\begin{equation}
\label{eq:ResKL}
\Res_{\mathrm{KL}}(P;\mathcal{R})
:=\inf_{Q\in\mathcal{R}}\KL(P\|Q).
\end{equation}
Escribimos $\Res_{\mathrm{ind}}:=\Res_{\mathrm{KL}}(P;\Rind)$ y
$\Res_{\mathrm{pair}}:=\Res_{\mathrm{KL}}(P;\Rpair)$.
\end{definition}

\begin{proposition}[Monoton\'ia al enriquecer la clase]
\label{prop:mono}
Si $\mathcal{R}\subset\mathcal{R}'$, entonces
$\Res_{\mathrm{KL}}(P;\mathcal{R}')\le\Res_{\mathrm{KL}}(P;\mathcal{R})$.
En particular
\begin{equation}
0\le\Res_{\mathrm{pair}}(P)\le\Res_{\mathrm{ind}}(P).
\end{equation}
\end{proposition}

\begin{proof}
El \'infimo sobre un conjunto m\'as grande no crece.
\end{proof}

\begin{remark}
El Nivel~3 \emph{fuerte} (irreductibilidad genuina a pares) es
$\Res_{\mathrm{pair}}>0$. El Nivel~3 \emph{d\'ebil} (no-factorizaci\'on total)
es $\Res_{\mathrm{ind}}>0$. El paradigma RECD apunta al fuerte; el proxy
$\excess$ mezcla se\~nales de ambos (Syn $\sim$ fuerte-heur\'istico;
Surp $\sim$ d\'ebil).
\end{remark}

% ============================================================
\section{Identidades exactas: TC, independencia, Surp}
\label{sec:identities}

\begin{theorem}[Total correlation como residuo a independencia]
\label{thm:TC}
Sea $P\in\mathcal{P}(\Sigma^d)$ con m\'argenes $P_i$ y
$P_{\mathrm{ind}}:=\prod_{i=1}^d P_i\in\Rind$. Entonces
\begin{equation}
\label{eq:TC_KL}
\TC(P)
:=\sum_{i=1}^d H(P_i)-H(P)
=\KL\bigl(P\,\|\,P_{\mathrm{ind}}\bigr)
=\Res_{\mathrm{ind}}(P).
\end{equation}
Adem\'as, el infimo que define $\Res_{\mathrm{ind}}$ se alcanza en
$P_{\mathrm{ind}}$ (producto de las m\'argenes \emph{de} $P$).
\end{theorem}

\begin{proof}
Para cualquier $Q=\prod q_i\in\Rind$,
\[
\KL(P\|Q)
=\sum_s P(s)\log\frac{P(s)}{\prod q_i(s_i)}
= -H(P)-\sum_i\sum_{a}P_i(a)\log q_i(a).
\]
La entrop\'ia cruzada $-\sum_a P_i(a)\log q_i(a)$ se minimiza en $q_i=P_i$,
con m\'inimo $H(P_i)$. Sumando,
$\inf_{Q\in\Rind}\KL(P\|Q)=\sum_i H(P_i)-H(P)=\TC(P)=\KL(P\|P_{\mathrm{ind}})$.
\end{proof}

\begin{corollary}
\label{cor:Res_ind_zero}
$\Res_{\mathrm{ind}}(P)=0$ sii $P=\prod P_i$ (independencia total).
\end{corollary}

\begin{definition}[Surprise de independencia del proxy]
\label{def:Surp}
Con $p_{\mathrm{obs}}=P$ y $p_{\mathrm{ind}}=P_{\mathrm{ind}}$,
\begin{equation}
\label{eq:Surp}
\Surp(P)
:=\sum_{s}P(s)\,
\Bigl[\log_2\frac{P(s)}{P_{\mathrm{ind}}(s)}\Bigr]_+.
\end{equation}
\end{definition}

\begin{proposition}[Relaci\'on Surp--TC]
\label{prop:Surp_TC}
\begin{enumerate}[label=(\roman*)]
  \item $\Surp(P)\ge 0$ y $\Surp(P)=0$ sii $P(s)\le P_{\mathrm{ind}}(s)$
  para todo $s$ en el soporte de $P$ (en particular si $P=P_{\mathrm{ind}}$).
  \item Descomponiendo el integrando de $\TC$ en partes positiva y
  negativa del information density
  $i(s):=\log_2\frac{P(s)}{P_{\mathrm{ind}}(s)}$,
  \begin{equation}
  \TC(P)=\mathbb{E}_P[i]_+-\mathbb{E}_P[i]_-
  =\Surp(P)-\mathbb{E}_P\bigl[(-i)_+\bigr],
  \end{equation}
  de modo que $\Surp(P)\ge\TC(P)$, con igualdad sii $i(s)\ge 0$ en todo el
  soporte de $P$.
\end{enumerate}
\end{proposition}

\begin{proof}
(i)~inmediato. (ii)~$\TC=\mathbb{E}_P[i]=\mathbb{E}_P[i]_+-\mathbb{E}_P[(-i)_+]$
y $\Surp=\mathbb{E}_P[i]_+$.
\end{proof}

\begin{remark}
Surp \emph{no} es un residuo a $\Rpair$: penaliza desviaciones por encima
del producto de m\'argenes, incluyendo estructura que un modelo de pares
podr\'ia explicar por completo.
\end{remark}

% ============================================================
\section{Syn, baseline de pares y residuo a $\Rpair$}
\label{sec:Syn}

\begin{definition}[Baseline de pares del proxy]
\label{def:Ipair}
Sea $\overline{I}:=\frac{2}{d(d-1)}\sum_{i<j}I(P_i;P_j)$ la MI media de
pares y
\begin{equation}
\label{eq:Ipair}
I_{\mathrm{pair}}(P):=(d-1)\,\overline{I}.
\end{equation}
(El factor $d-1$ es el contrato metodol\'ogico del paper de fundamentos /
methods excess3: hace $I_{\mathrm{pair}}$ dimensionalmente comparable a un
residual de multiinformaci\'on en un \'arbol o cadena de $d-1$ enlaces.)
\end{definition}

\begin{definition}[Syn del proxy]
\label{def:Syn}
\begin{equation}
\label{eq:Syn}
\Syn(P):=\bigl[\TC(P)-I_{\mathrm{pair}}(P)\bigr]_+.
\end{equation}
\end{definition}

\begin{proposition}[Propiedades elementales de Syn]
\label{prop:Syn_props}
\begin{enumerate}[label=(\roman*)]
  \item $0\le\Syn(P)\le\TC(P)=\Res_{\mathrm{ind}}(P)$.
  \item Si $d=2$, entonces $\TC(P)=I(P_1;P_2)$ e
  $I_{\mathrm{pair}}(P)=I(P_1;P_2)$, luego $\Syn(P)=0$ siempre.
  El Nivel~3 fuerte es vac\'io para $d=2$: hace falta $d\ge 3$.
  \item $\Syn$ \textbf{no} coincide en general con
  $\Res_{\mathrm{pair}}$: restar $(d-1)\overline{I}$ no es resolver el
  problema de proyecci\'on $\inf_{Q\in\Rpair}\KL(P\|Q)$.
\end{enumerate}
\end{proposition}

\begin{proof}
(i)--(ii)~c\'alculo directo. (iii)~$\Rpair$ maximiza entrop\'ia (minimiza
KL desde $P$) bajo constraints de m\'argenes bipolares; la soluci\'on es
el m\'aximo-entrop\'ia de orden 2, no la resta de un promedio de MI.
\end{proof}

\begin{definition}[Proyecci\'on de m\'axima entrop\'ia a pares]
\label{def:maxent2}
Sea $P^{(2)}\in\Rpair$ la (o una) proyecci\'on I-proyecci\'on de $P$ sobre
$\Rpair$: la ley de m\'axima entrop\'ia con las mismas bipolares
$\{P_{ij}\}$ que $P$ (cuando existe y es \'unica en el interior relativo).
Entonces
\begin{equation}
\label{eq:Res_pair_id}
\Res_{\mathrm{pair}}(P)=\KL\bigl(P\,\|\,P^{(2)}\bigr)
\end{equation}
bajo la existencia est\'andar de la I-proyecci\'on (alfabeto finito,
soporte compatible).
\end{definition}

\begin{theorem}[Descomposici\'on de la multiinformaci\'on]
\label{thm:decomp}
Siempre que exista $P^{(2)}$,
\begin{equation}
\label{eq:decomp}
\TC(P)
=\KL\bigl(P\,\|\,P_{\mathrm{ind}}\bigr)
=\KL\bigl(P\,\|\,P^{(2)}\bigr)
+\KL\bigl(P^{(2)}\,\|\,P_{\mathrm{ind}}\bigr),
\end{equation}
es decir
\begin{equation}
\Res_{\mathrm{ind}}(P)
=\Res_{\mathrm{pair}}(P)
+\TC\bigl(P^{(2)}\bigr).
\end{equation}
As\'i, $\TC$ se parte en: (a)~residuo irreducible a pares
$\Res_{\mathrm{pair}}$; (b)~multiinformaci\'on \emph{explicable por pares}
$\TC(P^{(2)})$.
\end{theorem}

\begin{proof}
Es la identidad de Pitman--Pythagoras para I-proyecciones en familias
exponenciales (o la cadena de KL cuando $P^{(2)}$ es la I-proyecci\'on de
$P$ sobre $\Rpair$ y $P_{\mathrm{ind}}$ es la I-proyecci\'on de $P^{(2)}$
sobre $\Rind\subset\Rpair$): 
$\KL(P\|P_{\mathrm{ind}})=\KL(P\|P^{(2)})+\KL(P^{(2)}\|P_{\mathrm{ind}})$.
\end{proof}

\begin{corollary}[Criterio limpio de sinergia irreducible fuerte]
\label{cor:strong}
$\Res_{\mathrm{pair}}(P)>0$ sii $P\ne P^{(2)}$: el joint no es un modelo de
pares. Equivalentemente, $\TC(P)>\TC(P^{(2)})$.
\end{corollary}

\begin{scholium}[Lectura para el paradigma]
La frase ``sinergia irreducible en t\'erminos ordinales'' se especializa
ahora a:
\[
\claim_3^{\mathrm{strong}}(t)
:\iff
\Res_{\mathrm{pair}}\bigl(\Jwin_t^{(w)}\bigr)>\theta_3.
\]
El Nivel~3 no es ``TC alto'': TC alto puede ser casi todo
$\TC(P^{(2)})$ (driver com\'un / red de pares).
\end{scholium}

% ============================================================
\section{Estatus exacto de $\excess$}
\label{sec:excess_status}

\begin{definition}[Contrato del proxy]
\label{def:excess}
\begin{equation}
\excess(P)=0.6\,\Syn(P)+0.4\,\Surp(P),
\end{equation}
con pesos fijos \emph{a priori} (nunca optimizados sobre el dataset
cient\'ifico bajo prueba).
\end{definition}

\begin{theorem}[Lo que $\excess$ \emph{s\'i} controla]
\label{thm:excess_yes}
\begin{enumerate}[label=(\roman*)]
  \item Si $P\in\Rind$, entonces $\TC=\Syn=\Surp=\excess=0$ y
  $\Res_{\mathrm{pair}}=0$.
  \item Si $\excess(P)>0$, entonces $P\notin\Rind$ (al menos una de
  $\Syn,\Surp$ es positiva $\Rightarrow$ desviaci\'on de independencia o
  residual de TC tras baseline de pares).
  \item Para $d=2$, $\Syn\equiv 0$ y $\excess=0.4\,\Surp$ colapsa a una
  funci\'on mon\'otona de la desviaci\'on positiva de independencia
  (no hay orden 3).
\end{enumerate}
\end{theorem}

\begin{proof}
(i)~Thm.~\ref{thm:TC} y Defs.~\ref{def:Syn}--\ref{def:Surp}.
(ii)~si ambas son 0: $\Surp=0$ y $\TC\le I_{\mathrm{pair}}$; si adem\'as
$P\in\Rind$ entonces todo es 0; si $\Surp=0$ y $\Syn=0$ a\'un podr\'ia
haber $P\notin\Rind$ con $i\le 0$ en el soporte y $\TC\le I_{\mathrm{pair}}$,
pero $\TC=\Res_{\mathrm{ind}}\ge 0$ y $\Surp\ge\TC$ (Prop.~\ref{prop:Surp_TC})
implican $\TC=0$, luego $P\in\Rind$. As\'i $\excess=0\Rightarrow P\in\Rind$
en esta cadena; contrapositiva: $\excess>0\Rightarrow P\notin\Rind$.
(iii)~Prop.~\ref{prop:Syn_props}(ii).
\end{proof}

\begin{theorem}[Lo que $\excess$ \emph{no} garantiza]
\label{thm:excess_no}
\begin{enumerate}[label=(\roman*)]
  \item \textbf{Falso:} $\excess>0\Rightarrow\Res_{\mathrm{pair}}>0$.
  Contraejemplo de clase: $P\in\Rpair\setminus\Rind$ (dependencia pura de
  pares). Entonces $\Res_{\mathrm{pair}}=0$, pero $\TC>0$ y t\'ipicamente
  $\Surp>0$, luego $\excess>0$ v\'ia Surp aunque no haya sinergia fuerte.
  \item \textbf{Falso:} $\Syn=0\Rightarrow\Res_{\mathrm{pair}}=0$.
  El baseline $(d-1)\overline{I}$ puede sobrestimar la masa de pares
  explicable, anulando Syn mientras $\KL(P\|P^{(2)})>0$.
  \item \textbf{Falso:} $\Syn>0\Rightarrow$ \'atomo sin\'ergico PID $>0$
  bajo toda descomposici\'on PID (depende del lattice y del estimador).
  \item Los pesos $0.6/0.4$ no son \'unicos ni \'optimos en ning\'un
  sentido de informaci\'on; son contrato de falsabilidad.
\end{enumerate}
\end{theorem}

\begin{example}[Driver com\'un vs candado XOR]
\label{ex:driver_xor}
\begin{itemize}
  \item \textbf{Driver com\'un:} $X_2,X_3$ copias ruidosas de $X_1$.
  Las bipolares capturan casi todo: $\Res_{\mathrm{pair}}\approx 0$,
  $\TC\approx\TC(P^{(2)})$, $\Syn\approx 0$, $\Surp$ puede ser alto
  $\Rightarrow$ $\excess$ inflado por Surp \emph{sin} Nivel~3 fuerte.
  \item \textbf{XOR / paridad:} $X_3=X_1\oplus X_2$ (alfabeto adecuado),
  pares marginalmente independientes. Entonces $\overline{I}\approx 0$,
  $\TC>0$, $\Syn\approx\TC$, $\Surp>0$, $\Res_{\mathrm{pair}}>0$:
  aqu\'i proxy y residuo fuerte co-direccionan.
\end{itemize}
\end{example}

\begin{scholium}[Diagn\'ostico operativo honesto]
Para no confundir driver com\'un con sinergia:
reportar \emph{separados} $\Syn$, $\Surp$, $\excess$, y cuando $d$ y
$|\Sigma|^d$ lo permitan, un estimador de $\Res_{\mathrm{pair}}$
(p.ej.\ $\KL(P\|P^{(2)})$ v\'ia iterative proportional fitting a
bipolares). El h\'ibrido $\excess$ es un indicador de screening, no el
veredicto de irreductibilidad.
\end{scholium}

% ============================================================
\section{Puente a PID y O-information}
\label{sec:PID}

\begin{definition}[O-information (recordatorio)]
\label{def:Oinfo}
La O-information de Rosas et al.\ es
\begin{equation}
\Omega(P)
=(d-2)\sum_i H(P_i)
+\sum_{i<j}H(P_{ij})
-\bigl[(d-1)H(P)+\cdots\bigr]
\end{equation}
(en la forma est\'andar para $d$ variables; signo positivo $\sim$
redundancia dominante, negativo $\sim$ sinergia dominante). No la
rederivamos: solo fijamos el rol.
\end{definition}

\begin{proposition}[Mapa de parentesco]
\label{prop:map}
\begin{center}
\begin{tabular}{@{}p{0.28\textwidth}p{0.32\textwidth}p{0.32\textwidth}@{}}
\toprule
Objeto & Captura & Relaci\'on con Nivel~3 \\
\midrule
$\TC=\Res_{\mathrm{ind}}$ & Toda dependencia vs independencia &
Necesario d\'ebil; no suficiente fuerte \\
$\Res_{\mathrm{pair}}$ & Residuo vs modelos de pares &
\textbf{Definici\'on fuerte de L3} \\
$\Syn$ & Heur\'istica $\TC-(d-1)\overline{I}$ &
Proxy de $\Res_{\mathrm{pair}}$, sesgable \\
$\Surp$ & Masa de co-ocurrencias sobre $P_{\mathrm{ind}}$ &
Se\~nal d\'ebil (L3 d\'ebil / pares) \\
$\excess$ & H\'ibrido fijo $0.6/0.4$ &
Screening preespecificado \\
O-information $\Omega$ & Balance sinergia/redundancia global &
Correlaciona con L3 en reg\'imenes; $\ne\Res$ \\
\'Atomos PID & Lattice nonneg.\ de partial info.\ &
Ideal local; no siempre estimable \\
\bottomrule
\end{tabular}
\end{center}
\end{proposition}

\begin{conjecture}[Puente cuantitativo Syn--$\Res_{\mathrm{pair}}$; refinada]
\label{conj:bridge}
Existe una subclase $\mathcal{P}_\ast\subset\mathcal{P}(\Sigma^d)$
(p.ej.\ $d\le 4$, $m=3$, $P$ de soporte lleno o acotado, ventanas $w$ con
concentraci\'on) y constantes $0<c\le C<\infty$ tales que
\begin{equation}
c\,\Res_{\mathrm{pair}}(P)
\le \Syn(P)
\le C\,\Res_{\mathrm{pair}}(P)
\qquad\text{para todo }P\in\mathcal{P}_\ast.
\end{equation}
\textbf{Estado:} abierta. La desigualdad izquierda es la delicada
(Syn puede anularse con $\Res_{\mathrm{pair}}>0$). La derecha es m\'as
plausible cuando $I_{\mathrm{pair}}$ no sobrecuenta en exceso.
Delimitaci\'on emp\'irica: validar en (i)~XOR, (ii)~driver com\'un,
(iii)~mapas log\'isticos acoplados pre/chaos del paper de fundamentos.
\end{conjecture}

\begin{conjecture}[Puente $\excess$--$\Res$ en la clase de trabajo RECD]
\label{conj:excess_bridge}
En la clase de leyes emp\'iricas generadas por mapas acoplados y por
sistemas con masa sin\'ergica dominante (no driver-puro),
$\excess$ es mon\'otono en media con $\Res_{\mathrm{pair}}$ y con la tasa
$\mathbb{P}(\Res_{\mathrm{pair}}>\theta)$. \textbf{No} se conjetura
equivalencia global (Thm.~\ref{thm:excess_no}).
\end{conjecture}

\begin{remark}[PID como extensi\'on natural A-IV-1]
Cuando un estimador PID estable est\'e disponible sobre $\Sigma^d$, la
sustituci\'on can\'onica es
\[
M_3^{\mathrm{PID}}(t):=S_{\mathrm{PID}}\bigl(\Jwin_t^{(w)}\bigr)
\]
(el \'atomo sin\'ergico del lattice elegido), manteniendo $\excess$ como
fallback de screening. Eso \emph{instancia} $\Res$ en un \'atomo, no
redefine el reloj.
\end{remark}

% ============================================================
\section{Condiciones de aparici\'on}
\label{sec:conditions}

\begin{proposition}[Condiciones necesarias]
\label{prop:nec}
Para $\Res_{\mathrm{pair}}(P)>0$ es necesario:
\begin{enumerate}[label=(\roman*)]
  \item $d\ge 3$;
  \item $P\notin\Rind$ (en particular $\TC(P)>0$);
  \item $P$ no es funci\'on de un grafo de pares exacto
  ($P\ne P^{(2)}$);
  \item el soporte de $P$ no se factoriza como producto de
  soportes bipolares consistentes con un MRF de pares sin ciclos de
  interacci\'on de orden 3 (en el sentido de Hammersley--Clifford para
  interacciones $\le 2$).
\end{enumerate}
\end{proposition}

\begin{proposition}[Condiciones suficientes (constructivas)]
\label{prop:suf}
Basta cualquiera de las siguientes:
\begin{enumerate}[label=(\roman*)]
  \item \textbf{Paridad / XOR generalizado:} existe una funci\'on
  $f:\Sigma^{d-1}\to\Sigma$ no representable por dependencias de pares
  tal que $P$ se concentra en el grafo de $x_d=f(x_1,\ldots,x_{d-1})$ y
  las bipolares son product\-like;
  \item \textbf{Interacci\'on log-lineal de orden 3:}
  $P\in$ familia exponencial con alg\'un $\theta_{ijk}\ne 0$ y
  proyecci\'on a $\Rpair$ que no reabsorbe esa interacci\'on;
  \item \textbf{Invariante que muere:} existe $I$ con $I\equiv 0$ en
  $\Rpair$ e $I(P)>0$ (Def.\ de Nota~1).
\end{enumerate}
\end{proposition}

\begin{conjecture}[Aparici\'on en la cascada de Feigenbaum; puente a pto.~4]
\label{conj:feigenbaum}
En familias de mapas unimodales acoplados d\'ebilmente, al cruzar la
acumulaci\'on de Feigenbaum hacia caos desarrollado,
$\mathbb{E}[\Res_{\mathrm{pair}}]$ y $\mathbb{E}[\Syn]$ crecen de forma no
explicada por un mero rescaling de $\mathbb{E}[|\tau_s|]$ o de
$\mathbb{E}[\Phi_1]$. Evidencia sint\'etica existente (exceso de mass
L3 pre$\to$chaos) es compatible; falta el teorema.
\end{conjecture}

% ============================================================
\section{Comportamiento bajo transformaciones}
\label{sec:transforms}

\begin{theorem}[Invariancia ordinal]
\label{thm:ordinal_inv}
Si $Z^{(i)}=f_i(X^{(i)})$ con cada $f_i$ estrictamente creciente, entonces
el proceso de s\'imbolos $S$ es id\'entico, y por tanto
$\TC$, $\Syn$, $\Surp$, $\excess$, $\Res_{\mathrm{ind}}$,
$\Res_{\mathrm{pair}}$ (sobre $S$) son id\'enticos.
\end{theorem}

\begin{proof}
Axioma A1 / Bandt--Pompe: las permutaciones no cambian.
\end{proof}

\begin{proposition}[Proyecci\'on a subconjuntos de variables]
\label{prop:marginalize}
Sea $P_A$ la marginal de $P$ sobre $A\subset\{1,\ldots,d\}$ con
$|A|\ge 3$. Entonces $\Res_{\mathrm{pair}}(P_A)$ puede ser positiva
aunque se calcule en un subconjunto; pero
$\Res_{\mathrm{pair}}(P)=0$ \emph{no} implica $\Res_{\mathrm{pair}}(P_A)=0$
ni rec\'iprocamente de forma libre (la irreductibilidad no es mon\'otona
trivial bajo marginaci\'on).
\end{proposition}

\begin{proposition}[Qu\'e se pierde al proyectar a pares]
\label{prop:loses}
La aplicaci\'on $P\mapsto P^{(2)}$ (o $P\mapsto\{P_{ij}\}$) anula todo
funcional que sea invariante de orden superior. En particular
$\Res_{\mathrm{pair}}(P^{(2)})=0$ y todo \'atomo sin\'ergico PID de $P^{(2)}$
es nulo. Eso es la formulaci\'on operativa de ``invariantes que
desaparecen al proyectar a niveles inferiores''.
\end{proposition}

\begin{proposition}[Permutaci\'on de canales]
\label{prop:perm}
Si $\sigma\in S_d$ reordena coordenadas,
$\Res_{\mathrm{pair}}(P\circ\sigma^{-1})=\Res_{\mathrm{pair}}(P)$ y lo
mismo para $\TC,\Syn,\Surp,\excess$ (simetr\'ia del joint).
\end{proposition}

% ============================================================
\section{Protocolo de uso te\'orico-operativo}
\label{sec:protocol}

\begin{enumerate}[label=\textbf{P\arabic*.}]
  \item \textbf{Definir} Nivel~3 fuerte como $\Res_{\mathrm{pair}}>\theta_3$
  (claim).
  \item \textbf{Estimar} cuando sea factible: $P^{(2)}$ por IPF/maxent
  de pares y $\KL(P\|P^{(2)})$; si no, usar $\Syn$ como cota/proxy
  principal de irreductibilidad y $\Surp$ como se\~nal de
  no-independencia.
  \item \textbf{Reportar} siempre el triple $(\Syn,\Surp,\excess)$
  separado, nunca solo el escalar h\'ibrido.
  \item \textbf{Contrastar} con nulo de independencia y, si es posible,
  con surrogados que preservan bipolares (destruyen orden $\ge 3$).
  \item \textbf{No afirmar} PID completo, causalidad, ni equivalencia
  $\excess\Leftrightarrow\Res_{\mathrm{pair}}$.
  \item \textbf{RECD:} $\widehat{M}_3=\excess$ sigue siendo admisible
  como evidencia de screening en $\Delta\RECD$; la interpretaci\'on de
  ``masa ontol\'ogica de L3'' debe leerse como masa de
  \emph{proxy de residual}, no como masa del \'atomo PID.
\end{enumerate}

% ============================================================
\section{Respuestas al criterio de \'exito (punto 2)}
\label{sec:answers}

\begin{enumerate}[label=\textbf{Q\arabic*.}, leftmargin=*]
  \item \textbf{¿Qu\'e es sinergia irreducible en t\'erminos ordinales?}\\
  $\Res_{\mathrm{pair}}(\Jwin_t^{(w)})=\KL(\Jwin\|P^{(2)})>\theta_3$:
  el joint de s\'imbolos no es un modelo de interacciones de orden
  $\le 2$.

  \item \textbf{¿Por qu\'e el Nivel~3 no es ``m\'as correlaci\'on''?}\\
  Correlaci\'on / $\tau_s$ / $\Phi_1$ / $\Phi_2$ / $\TC(P^{(2)})$ viven
  en $\Rpair$ o en observables de orden $\le 2$. L3 es el sumando
  $\KL(P\|P^{(2)})$ de la descomposici\'on
  $\TC(P)=\Res_{\mathrm{pair}}+\TC(P^{(2)})$ (Thm.~\ref{thm:decomp}).

  \item \textbf{¿Qu\'e aporta $\excess$?}\\
  Un proxy preespecificado y computable que (a)~se anula en independencia
  total, (b)~detecta geometr\'ias tipo XOR v\'ia Syn, (c)~mantiene visible
  la co-ocurrencia via Surp, (d)~no certifica $\Res_{\mathrm{pair}}>0$ cuando
  Surp domina por pares.

  \item \textbf{¿Condiciones de aparici\'on?}\\
  Necesarias: $d\ge 3$, no-factorizaci\'on a pares.
  Suficientes constructivas: XOR/paridad, interacci\'on log-lineal de
  orden 3, invariante que muere en $\Rpair$.
  En din\'amica (Feigenbaum): conjetura abierta (pto.~4).
\end{enumerate}

% ============================================================
\section{Canon a\~nadido (C8--C14)}
\label{sec:canon}

\begin{enumerate}[label=\textbf{C\arabic*.}]
  \setcounter{enumi}{7}
  \item $\Res_{\mathrm{ind}}(P)=\TC(P)=\KL(P\|P_{\mathrm{ind}})$.
  \item $\TC(P)=\Res_{\mathrm{pair}}(P)+\TC(P^{(2)})$ (descomposici\'on).
  \item L3 \emph{fuerte} $:=$ $\Res_{\mathrm{pair}}>0$; L3 \emph{d\'ebil}
  $:=$ $\Res_{\mathrm{ind}}>0$.
  \item $\Syn$ es heur\'istica, no I-proyecci\'on a $\Rpair$.
  \item $\Surp\ge\TC$, con igualdad solo si el information density no
  cambia de signo; Surp no separa pares de orden 3.
  \item $\excess>0\Rightarrow$ L3 d\'ebil; $\not\Rightarrow$ L3 fuerte.
  \item Reportar $(\Syn,\Surp,\excess)$ por separado; surrogados que
  preservan bipolares para testear orden $\ge 3$.
\end{enumerate}

\vspace{1.5em}
\noindent\rule{\textwidth}{0.4pt}\\[0.5em]
\begin{center}
\em Fin de la Nota formal punto 2 v0.1.\\
Pr\'oximo entregable natural: estatus de $\alpha(\lambda)$ (punto 3),
o validaci\'on emp\'irica de Conjs.~\ref{conj:bridge}--\ref{conj:feigenbaum}.
\end{center}

\begin{thebibliography}{99}
\bibitem{Padilla2026note1}
J.~Padilla-Villanueva.
\newblock Nota formal: Tres niveles de conjunci\'on ordinal\ldots\ (1a--1c).
\newblock Fundaci\'on RECD, 2026.
\bibitem{Padilla2026framework}
J.~Padilla-Villanueva.
\newblock Systemic Tau and the RECD Framework.
\newblock Manuscrito de fundamentos, 2026.
\bibitem{Watanabe1960}
S.~Watanabe.
\newblock Information theoretical analysis of multivariate correlation.
\newblock \emph{IBM J.\ Res.\ Dev.} 4, 66--82 (1960).
\bibitem{Williams2010}
P.~L.~Williams and R.~D.~Beer.
\newblock Nonnegative decomposition of multivariate information.
\newblock arXiv:1004.2515 (2010).
\bibitem{Rosas2019}
F.~E.~Rosas et~al.
\newblock Quantifying high-order interdependencies via multivariate
extensions of the mutual information.
\newblock \emph{Phys.\ Rev.\ E} 100, 032305 (2019).
\bibitem{Csiszar}
I.~Csisz\'ar and F.~Mat\'u\v{s}.
\newblock Information projections revisited.
\newblock \emph{IEEE Trans.\ Inf.\ Theory} 49, 1474--1490 (2003).
\end{thebibliography}

\end{document}
